% Generated draft: requires mathematical and source review.
\documentclass[11pt]{article}
\usepackage[margin=25mm]{geometry}
\usepackage{fontspec}
\setmainfont{Latin Modern Roman}
\usepackage{amsmath,amssymb,mathtools,unicode-math}
\setmathfont{Latin Modern Math}
\usepackage{enumitem,needspace,xurl,hyperref}
\setlength{\parindent}{0pt}
\setlength{\parskip}{5pt}
\setlength{\emergencystretch}{3em}
\newcommand{\N}{\mathbb{N}}
\newcommand{\Z}{\mathbb{Z}}
\newcommand{\Q}{\mathbb{Q}}
\newcommand{\R}{\mathbb{R}}
\newcommand{\C}{\mathbb{C}}
\newcommand{\F}{\mathbb{F}}
\newcommand{\A}{\mathbb{A}}
\newcommand{\PP}{\mathbb{P}}
\newcommand{\E}{\mathbb{E}}
\newcommand{\eps}{\varepsilon}
\newcommand{\dd}{\,\mathrm{d}}
\begin{document}
\section*{UnsolvedMath ID: 701351}
\textbf{Source problem code:} catalogue entry
\subsection{Definitions and mathematical statement}
A topological group is a group \(G\) equipped with a topology \(\tau\) for which multiplication \(G\times G\to G\) and inversion \(G\to G\) are continuous, with the product topology on \(G\times G\). The supplied source does not specify a separation convention, such as requiring Hausdorffness. In the statements below, the quantifier over topological groups must use the convention intended by the source; this qualification remains unresolved.

The topology is non-discrete if \(\tau\ne\mathcal P(G)\). A subset \(A\subseteq G\) is nowhere dense if
\[
\operatorname{int}_G\bigl(\overline A^{\,G}\bigr)=\varnothing.
\]
A point \(a\in A\) is isolated in \(A\) if there is an open set \(U\in\tau\) such that \(U\cap A=\{a\}\). Thus \(A\) has no isolated points if
\[
\forall a\in A\ \forall U\in\tau,\qquad
 a\in U\ \Longrightarrow\ (U\cap A)\setminus\{a\}\ne\varnothing.
\]

\textbf{Consistency question.} Let \(\Phi\) be the assertion that every non-discrete topological group \((G,\tau)\) contains a subset \(A\subseteq G\) satisfying
\[
A\ne\varnothing,\qquad
\operatorname{int}_G\bigl(\overline A^{\,G}\bigr)=\varnothing,
\qquad A\text{ has no isolated points}.
\]
Is \(\mathrm{ZFC}+\Phi\) consistent?

\textbf{Supplied additional assertion.} The record states that, under Martin's Axiom, there exists a countable non-discrete topological group \((G,\tau)\) such that
\[
\forall A\subseteq G,\quad
\left(A\ne\varnothing\ \land\
\operatorname{int}_G\bigl(\overline A^{\,G}\bigr)=\varnothing\right)
\Longrightarrow
\exists a\in A\ \exists U\in\tau\quad U\cap A=\{a\}.
\]
This additional assertion is reproduced from the supplied record, without independent verification. The record does not specify a version of Martin's Axiom or any additional set-theoretic hypotheses.
\subsection{Short English statement}
Is it consistent with ZFC that every non-discrete topological group has a nonempty nowhere dense subset with no isolated points in its subspace topology? The supplied record also asserts that Martin's Axiom yields countable counterexamples to this universal property. The intended separation convention for topological groups and the precise Martin's Axiom hypothesis require source clarification.
\subsection{Sources}
\begin{itemize}
\item Local UnsolvedMath registry, supplied complete record: permanent ID \(701351\), exact source problem code \texttt{catalogue entry}, title \emph{Kourovka Notebook Problem 17.49}. The registry labels the problem \texttt{open}; this is supplied status metadata, not an independently verified assessment.
\item E. G. Zelenyuk, proposer, 2010. The supplied background attributes Problem 17.49 to \emph{The Kourovka Notebook (Unsolved Problems in Group Theory)}, 17th issue (2010), and references the 21st issue edition (2026), \texttt{arXiv:1401.0300v48}.
\item The supplied collection description names E. I. Khukhro and V. D. Mazurov as editors of the Kourovka Notebook collection covering Issues 1--21 (1965--2026). These bibliographic details and the assertion concerning Martin's Axiom are supplied attribution and have not been independently checked.
\end{itemize}
\end{document}
